% =============================================================
% Kapitel 6: Schlussbetrachtung
% =============================================================
\chapter{Schlussbetrachtung}
\label{chapter:kapitel6}

Ziel der vorliegenden Arbeit war es zu bestimmen, inwieweit die bei Woodward L'Orange in der mobilen Instandhaltung eingesetzte SAP-Fiori-Anwendung zur Fotodokumentation die Interaktionsprinzipien der ISO~9241-110:2020 erfüllt und welche Optimierungspotenziale sich aus einer heuristischen Evaluation ableiten lassen.
Dazu wurde ein Kriterienkatalog entwickelt, der die Interaktionsprinzipien der Norm in überprüfbare Prüffragen überführt, und die Anwendung im Rahmen einer Experteninspektion systematisch gegen diesen Katalog geprüft.

Die Evaluation ergab ein heterogenes Bild über die sieben Interaktionsprinzipien hinweg.
Die Prinzipien Aufgabenangemessenheit, Erlernbarkeit, Steuerbarkeit und Benutzerbindung wurden als erfüllt bewertet, die Selbstbeschreibungsfähigkeit und die Erwartungskonformität als teilweise erfüllt.
Als zentraler Schwachpunkt erwies sich die Robustheit gegen Benutzungsfehler, die nicht erfüllt wurde, da beim Hochladen ungültiger Dateiformate ein Fehlerfeedback fehlt und die bei Ladefehlern erscheinenden Meldungen weder den eigentlichen Fehler benennen noch handlungsleitend formuliert sind.
Aus den insgesamt elf identifizierten Abweichungen wurden nach Schweregrad priorisierte Optimierungspotenziale abgeleitet.
Die Forschungsfrage lässt sich somit dahingehend beantworten, dass die Anwendung die Interaktionsprinzipien der ISO~9241-110:2020 überwiegend, jedoch nicht vollständig erfüllt, wobei die praktisch relevanten Defizite vor allem im Umgang mit Benutzungsfehlern liegen.

Die Befunde sind vor dem Hintergrund der methodischen Grenzen der Untersuchung einzuordnen.
An der ursprünglichen Entwicklung der Anwendung war der Verfasser nicht beteiligt und kannte diese zum Zeitpunkt der Evaluation noch nicht, sodass die Bewertung unabhängig von den ursprünglichen Entwicklungsentscheidungen erfolgte.
Da er die Abweichungen im Anschluss jedoch selbst als Entwickler behebt, lässt sich ein unbewusster Einfluss dieser Zuständigkeit auf die Schweregradbeurteilung nicht vollständig ausschließen.
Hinzu kommt, dass eine heuristische Evaluation durch eine einzelne Person nur einen Teil der tatsächlich vorhandenen Usability-Probleme aufdeckt.\footcite[Vgl.][S.~209]{nielsen1993mathematical}
Die Ergebnisse sind daher als Momentaufnahme zu verstehen.

Trotz dieser Grenzen bietet der entwickelte Kriterienkatalog einen Mehrwert über die konkrete Untersuchung hinaus, da er grundsätzlich auf weitere mobile SAP-Fiori-Anwendungen im Instandhaltungskontext übertragbar ist.
Für weiterführende Arbeiten bietet es sich an, die heuristischen Befunde durch Nutzertests mit den Instandhaltungstechnikern empirisch zu validieren sowie eine kollaborative Inspektion mit mehreren, insbesondere anwendungsfremden Evaluatoren durchzuführen, um eine breitere Problemabdeckung zu erzielen.\footcite[Vgl.][S.~363]{inal2021collaborative}
Für Woodward L'Orange ergibt sich als konkreter nächster Schritt die Behebung der Schwachstellen bei der Robustheit gegen Benutzungsfehler, insbesondere die Ergänzung eines Fehlerfeedbacks bei ungültigen Datei-Uploads und die Überarbeitung der Fehlermeldungen hin zu handlungsleitenden Hinweisen.
